% 可数集（综述）
% license CCBYSA3
% type Wiki

本文根据 CC-BY-SA 协议转载翻译自维基百科\href{https://en.wikipedia.org/wiki/Countable_set}{相关文章}

在数学中，如果一个集合是有限的，或者它可以与自然数集合建立一一对应关系，那么它就是可数的。[a] 等价地说，如果存在一个从该集合到自然数的单射函数，那么该集合就是可数的；这意味着集合中的每个元素都可以与一个唯一的自然数对应，或者说集合的元素可以一个接一个地数出来，尽管由于元素个数无限，计数过程可能永远不会结束。

从更技术性的角度讲，假设接受可数选择公理，若一个集合的基数（即集合中的元素个数）不大于自然数集合的基数，则它是可数的。一个可数但不是有限的集合称为可数无穷。

这一概念归功于格奥尔格·康托尔，他证明了不可数集的存在，也就是说存在一些集合不是可数的；例如实数集合就是不可数的。
\subsection{关于术语的说明}
尽管这里所定义的“可数”与“可数无穷”用法相当普遍，但并非所有文献都采用这种术语。[1] 一种替代用法是：用countable 来表示这里所说的“可数无穷”，而用at most countable来表示这里所说的“可数”。[2][3]

术语enumerable[4]和denumerable[5][6] 也可能被使用，例如分别指代“可数”和“可数无穷”。[7] 但其定义并不统一，并且需要注意与“递归可枚举”的区别。[8]
\subsection{定义}
一个集合 $S$ 是**可数的**，当且仅当满足以下任一条件：
\begin{itemize}
\item * 它的基数 $|S|$ 小于或等于 $\aleph_{0}$（阿列夫零，$\mathbb{N}$ ——自然数集的基数）。\[9]
* 存在一个从 $S$ 到 $\mathbb{N}$ 的单射函数。\[10]\[11]
* $S$ 是空集，或者存在一个从 $\mathbb{N}$ 到 $S$ 的满射函数。\[11]
* 存在 $S$ 与 $\mathbb{N}$ 的某个子集之间的双射映射。\[12]
* $S$ 要么是有限的（$|S| < \aleph_{0}$），要么是可数无穷的。\[5]
\end{itemize}
以上定义都是**等价的**。

一个集合 $S$ 是**可数无穷的**，当且仅当：
\begin{itemize}
\item 它的基数 $|S|$ 正好等于 $\aleph_{0}$。[9]
\item 存在一个从 $S$ 到 $\mathbb{N}$ 的既单射又满射（因此是双射）的映射。
\item $S$ 与 $\mathbb{N}$ 之间存在一一对应关系。[13]
\item $S$ 的元素可以排列成一个无限序列$a_{0}, a_{1}, a_{2}, \ldots$
其中 $a_i \neq a_j$（当 $i \neq j$ 时），并且 $S$ 的每个元素都在该序列中出现。[14][15]
\end{itemize}
如果一个集合不是可数的，即它的基数大于 $\aleph_{0}$，那么它就是不可数的。[9]

